\section{Dataset Description} \label{sec:dataset}

\subsection{Data Collection and Composition}

We constructed a plagiarism detection dataset comprising 100+ text pairs derived from AI and machine learning research paper summaries. The dataset captures diverse writing styles, scientific terminology, and subject matter variation to ensure robust model evaluation.

\subsubsection{Source Material}

Data sources include:
\begin{itemize}
  \item ArXiv abstracts (machine learning, NLP, computer vision papers)
  \item IEEE Xplore conference proceedings
  \item ACM Digital Library publications
  \item Academic journal articles in computer science
\end{itemize}

This selection ensures domain-specific vocabulary, technical complexity, and authentic academic writing patterns representative of institutional plagiarism detection use cases.

\subsection{Plagiarism Variant Types}

The dataset includes balanced representation of plagiarism types:

\begin{table}[h!]
\centering
\caption{Dataset Composition by Plagiarism Type}
\label{tab:dataset_composition}
\begin{tabular}{|l|c|c|}
\hline
\textbf{Plagiarism Category} & \textbf{Count} & \textbf{Percentage} \\
\hline
Direct Copy & 25 & 25\% \\
Paraphrasing & 25 & 25\% \\
Semantic Variation & 20 & 20\% \\
Non-plagiarized (Original) & 30 & 30\% \\
\hline
\textbf{Total} & \textbf{100} & \textbf{100\%} \\
\hline
\end{tabular}
\end{table}

\subsubsection{Plagiarism Type Definitions}

\begin{itemize}
  \item \textbf{Direct Copy} --- Verbatim copying with $\geq 90\%$ lexical overlap
  \item \textbf{Paraphrasing} --- Semantic equivalence with 30-70\% lexical overlap through synonym substitution and structural rearrangement
  \item \textbf{Semantic Variation} --- Content expressing similar ideas (10-30\% lexical overlap) with different terminology and examples
  \item \textbf{Non-plagiarized} --- Completely original content with $\leq 10\%$ overlap to legitimate common knowledge
\end{itemize}

\subsection{Data Preprocessing Pipeline}

All input text undergoes standardized preprocessing:

\begin{enumerate}
  \item \textbf{Lowercasing} --- Convert all text to lowercase for uniform representation
  
  \item \textbf{URL/Email Removal} --- Strip web links and email addresses not relevant to content similarity
  
  \item \textbf{Special Character Filtering} --- Remove punctuation except hyphens in compound words
  
  \item \textbf{Whitespace Normalization} --- Collapse multiple spaces; convert tabs to spaces
  
  \item \textbf{Tokenization} --- Split into word tokens for model ingestion
\end{enumerate}

No stopword removal is performed, as these often carry semantic significance in academic writing (e.g., ``with'', ``without'' in methodology descriptions).

\subsection{Data Splits and Balance}

\begin{table}[h!]
\centering
\caption{Dataset Partitioning}
\label{tab:splits}
\begin{tabular}{|l|c|c|c|}
\hline
\textbf{Split} & \textbf{Pairs} & \textbf{Plagiarized} & \textbf{Original} \\
\hline
Training & 70 & 35 & 35 \\
Validation & 15 & 7 & 8 \\
Test & 15 & 8 & 7 \\
\hline
\end{tabular}
\end{table}

Stratified splitting ensures balanced class distribution (approximately 50\% plagiarized, 50\% original) across all partitions, preventing training bias toward either class.

\subsection{Dataset Characteristics}

\begin{itemize}
  \item \textbf{Average text length} --- 150-300 tokens per document pair
  \item \textbf{Vocabulary diversity} --- 2,500+ unique terms including domain-specific terminology
  \item \textbf{Linguistic complexity} --- Academic writing with complex syntax and technical vocabulary
  \item \textbf{Temporal diversity} --- Papers spanning 2015-2023, capturing evolving terminology
\end{itemize}

These characteristics ensure the dataset reflects real-world institutional plagiarism detection scenarios rather than simplified toy problems.
